% Источники: exam/materials/моделирование.pdf, раздел 30; exam/materials/преподаватель/2020/23-03-2020-MODEL L6(2).doc.
\section{Анализ функциональных зависимостей между наблюдаемыми одновременно характеристиками объекта. Метод наименьших квадратов.}

\subsection*{Постановка}

Требуется подобрать параметры $\boldsymbol{\beta}$ функции $f(X,\boldsymbol{\beta})$ так, чтобы она наилучшим образом описывала наблюдения $(X_t, Y_t)$, $t=1,\ldots,n$.

\subsection*{Критерий МНК}

Минимизируется сумма квадратов отклонений:
\[
  F(\boldsymbol{\beta}) = \sum_{t=1}^{n} \bigl(Y_t - f(X_t,\boldsymbol{\beta})\bigr)^2 \to \min_{\boldsymbol{\beta}}.
\]

\subsection*{Линейный МНК (простая регрессия)}

Модель $Y_t = \beta_0 + \beta_1 X_t + \varepsilon_t$. Нормальные уравнения:
\[
  \begin{pmatrix} n & \sum x_t \\ \sum x_t & \sum x_t^2 \end{pmatrix}
  \begin{pmatrix} \hat{\beta}_0 \\ \hat{\beta}_1 \end{pmatrix}
  = \begin{pmatrix} \sum y_t \\ \sum x_t y_t \end{pmatrix}.
\]
Решение:
\[
  \hat{b} = \frac{\sum(x_t-\bar{x})(y_t-\bar{y})}{\sum(x_t-\bar{x})^2}, \qquad \hat{a} = \bar{y} - \hat{b}\bar{x}.
\]

\subsection*{Множественная линейная регрессия}

Модель: $\mathbf{Y} = X\boldsymbol{\beta} + \boldsymbol{\varepsilon}$, где $X$ — матрица плана ($n \times p$).

МНК-оценка:
\[
  \hat{\boldsymbol{\beta}} = (X^\top X)^{-1} X^\top \mathbf{Y}.
\]
\textbf{Теорема Гаусса–Маркова:} при $M[\boldsymbol{\varepsilon}]=\mathbf{0}$ и $\mathrm{Cov}(\boldsymbol{\varepsilon})=\sigma^2 I$ МНК-оценка $\hat{\boldsymbol{\beta}}$ является лучшей линейной несмещённой оценкой (BLUE).

\subsection*{Оценка качества подгонки}

Коэффициент детерминации:
\[
  R^2 = 1 - \frac{\sum(y_t - \hat{y}_t)^2}{\sum(y_t - \bar{y})^2} \in [0,1].
\]
$R^2 = 1$ — идеальная подгонка; $R^2 = 0$ — модель не лучше среднего.

\clearpage
